% =============================================================================
% test-ibge-tabelas.tex
% Documento de teste para ibge-tabelas.sty
%
% Demonstra formatação de tabelas conforme Normas de Apresentação Tabular
% do IBGE (3ª edição, 1993).
%
% Compilação:
%   TEXINPUTS=./src: latexmk -pdf -output-directory=tests tests/test-ibge-tabelas.tex
%
% Checklist:
%   1  — Três traços horizontais (4.3.1)
%   2  — Sem bordas verticais (4.3.3)
%   3  — Traços verticais internos opcionais (4.3.2)
%   4  — Numeração automática (4.1)
%   5  — Título no topo (4.2)
%   6  — Fonte obrigatória no rodapé (4.10)
%   7  — Ordem do rodapé: Fonte → Nota → Notas específicas (4.10–4.12)
%   8  — Sinais convencionais (4.8.1)
%   9  — Chamada remissiva (4.9)
%   10 — Vírgula decimal (4.7)
%   11 — Classes de frequência (6)
%   12 — Apresentação de tempo (5)
%   13 — Nota de arredondamento (7)
%   14 — Compatibilidade hyperref
%   15 — Warning sem fonte (4.10)
%   16 — Rodapé avulso com \abntRodape (4.10)
%   17 — Chamada em colchetes / exponencial (4.9.1)
%   18 — Numeração subordinada a capítulos (4.1.2)
% =============================================================================

\documentclass[a4paper, 12pt]{report}
\usepackage[T1]{fontenc}
\usepackage[brazil]{babel}
\usepackage{abnt-numeracao}
\usepackage{ibge-tabelas}
\usepackage{hyperref}                % checklist 14
\usepackage{setspace}
\onehalfspacing

\begin{document}

% =====================================================================
% CAPÍTULO 1 — Tabela básica (checklist 1, 2, 4, 5, 6, 10)
% =====================================================================
\chapter{Tabelas básicas}

A Tabela~\ref{tab:producao} apresenta dados de produção por estado.

% Checklist 1 (três traços), 2 (sem bordas verticais), 4 (numeração),
% 5 (título no topo), 6 (fonte no rodapé), 10 (vírgula decimal)
\begin{tabela}{Produção de casulos do bicho-da-seda, por Unidade da
Federação \textendash\ Brasil \textendash\ 1974}
  \label{tab:producao}
  \abntFonte{IBGE, Diretoria de Pesquisas, Departamento de Agropecuária.}
  \begin{tabelaabnt}{colspec={lrr}}
    Unidade da Federação & Produção (t) & Participação (\%) \\
    São Paulo            & \num{120.50} & \num{45.30}       \\
    Paraná               & \num{98.20}  & \num{36.90}       \\
    Minas Gerais         & \num{32.80}  & \num{12.30}       \\
    Outros               & \num{14.50}  & \num{5.50}        \\
    Total                & \num{266.00} & \num{100.00}      \\
  \end{tabelaabnt}
\end{tabela}

% =====================================================================
% CAPÍTULO 2 — Rodapé completo e chamadas (checklist 7, 9)
% =====================================================================
\chapter{Rodapé e chamadas}

A Tabela~\ref{tab:populacao} demonstra o rodapé completo.

% Checklist 7 (ordem do rodapé), 9 (chamada)
\begin{tabela}{População residente, por sexo \textendash\ Brasil
\textendash\ 2022}
  \label{tab:populacao}
  \abntFonte{IBGE, Censo Demográfico 2022.}
  \abntNota{Dados sujeitos a retificação.}
  \abntNotaEsp{1}{Inclui nascidos vivos registrados em anos anteriores}
  \abntNotaEsp{2}{Exclusive a população rural de Rondônia, Acre, Amazonas,
  Roraima, Pará e Amapá}
  \begin{tabelaabnt}{colspec={lrrr}}
    Região     & Homens\abntChamada{1} & Mulheres & Total\abntChamada{2} \\
    Norte      & \num{9200000}   & \num{9100000}  & \num{18300000}  \\
    Nordeste   & \num{27500000}  & \num{29000000} & \num{56500000}  \\
    Sudeste    & \num{42100000}  & \num{44900000} & \num{87000000}  \\
    Sul        & \num{14800000}  & \num{15200000} & \num{30000000}  \\
    Centro-Oeste & \num{8200000} & \num{8300000}  & \num{16500000}  \\
  \end{tabelaabnt}
\end{tabela}

% =====================================================================
% CAPÍTULO 3 — Sinais convencionais (checklist 8)
% =====================================================================
\chapter{Sinais convencionais}

A Tabela~\ref{tab:sinais} demonstra os sinais convencionais da norma.

\begin{tabela}{Demonstração dos sinais convencionais}
  \label{tab:sinais}
  \abntFonte{Dados fictícios para fins de demonstração.}
  \abntNota{Os sinais convencionais utilizados seguem as Normas de
  Apresentação Tabular do IBGE.}
  \begin{tabelaabnt}{colspec={lrrr}}
    Indicador     & 2020            & 2021            & 2022            \\
    Produção A    & \num{1500.20}   & \num{1620.50}   & \num{1780.30}   \\
    Produção B    & \abntSinalZero  & \num{42.10}     & \num{85.60}     \\
    Produção C    & \abntSinalNA    & \abntSinalNA    & \num{10.00}     \\
    Produção D    & \abntSinalND    & \num{320.00}    & \num{340.00}    \\
    Produção E    & \abntSinalOmitido & \abntSinalOmitido & \num{5.00}  \\
    Arredondado+  & \abntSinalArrZero[2] & \num{0.10} & \num{0.20}     \\
    Arredondado$-$ & \abntSinalArrZero*[1] & \num{-0.3} & \num{0.10}  \\
  \end{tabelaabnt}
\end{tabela}

% =====================================================================
% Texto com classes de frequência e apresentação de tempo
% (checklist 11, 12, 13)
% =====================================================================
\chapter{Exemplos auxiliares}

\section{Apresentação de tempo}

% Checklist 12
A série de dados abrange o período de \abntSerieConsec{1981}{1985}
(série consecutiva). Para anos não consecutivos, usa-se
\abntSerieNConsec{1981}{1985}.

\section{Classes de frequência}

% Checklist 11
As classes de frequência podem ser expressas de três formas:
\begin{itemize}
  \item Inclui inferior, exclui superior: \abntClasseIE{10}{20}
  \item Exclui inferior, inclui superior: \abntClasseEI{10}{20}
  \item Inclui ambos: \abntClasseII{10}{20}
\end{itemize}

\section{Nota de arredondamento}

% Checklist 13
\abntNotaArrend

% =====================================================================
% Tabela sem fonte — gera warning (checklist 15)
% =====================================================================
\begin{tabela}{Tabela sem fonte para teste de warning}
  \label{tab:semfonte}
  \begin{tabelaabnt}{colspec={lr}}
    Item & Valor \\
    A    & \num{10} \\
  \end{tabelaabnt}
\end{tabela}

% =====================================================================
% CAPÍTULO 5 — Rodapé avulso (checklist 16)
% =====================================================================
\chapter{Rodapé avulso}

A Tabela~\ref{tab:avulso} demonstra o uso de \verb|\abntRodape| fora do
ambiente \texttt{tabela}, imprimindo o rodapé de forma avulsa.

\begin{tabela}{Dados com rodapé impresso manualmente}
  \label{tab:avulso}
  \abntFonte{Elaboração própria.}
  \abntNota{Rodapé será impresso fora do ambiente tabela.}
  \begin{tabelaabnt}{colspec={lr}}
    Indicador & Valor \\
    X         & \num{42} \\
    Y         & \num{58} \\
  \end{tabelaabnt}
\end{tabela}

% Checklist 16 — uso avulso do rodapé
\abntRodape

% =====================================================================
% CAPÍTULO 6 — Chamada em colchetes e exponencial (checklist 17)
% =====================================================================
\chapter{Formatos de chamada}

\section{Chamada em colchetes}

\begingroup
\ExplSyntaxOn \keys_set:nn { ibge-tabelas } { chamada=colchetes } \ExplSyntaxOff
\begin{tabela}{Tabela com chamada em colchetes}
  \label{tab:colchetes}
  \abntFonte{Dados fictícios.}
  \abntNotaEsp{1}{Nota com chamada em colchetes}
  \begin{tabelaabnt}{colspec={lr}}
    Item\abntChamada{1} & Valor \\
    A                   & \num{100} \\
  \end{tabelaabnt}
\end{tabela}
\endgroup

\section{Chamada exponencial}

\begingroup
\ExplSyntaxOn \keys_set:nn { ibge-tabelas } { chamada=exponencial } \ExplSyntaxOff
\begin{tabela}{Tabela com chamada exponencial}
  \label{tab:exponencial}
  \abntFonte{Dados fictícios.}
  \abntNotaEsp{1}{Nota com chamada exponencial}
  \begin{tabelaabnt}{colspec={lr}}
    Item\abntChamada{1} & Valor \\
    B                   & \num{200} \\
  \end{tabelaabnt}
\end{tabela}
\endgroup

% =====================================================================
% CAPÍTULO 7 — Numeração subordinada (checklist 18)
% =====================================================================
\chapter{Numeração subordinada}

\begingroup
\ExplSyntaxOn \keys_set:nn { ibge-tabelas } { subordinar=true } \ExplSyntaxOff
\begin{tabela}{Tabela com numeração subordinada ao capítulo}
  \label{tab:subordinada}
  \abntFonte{Dados fictícios.}
  \begin{tabelaabnt}{colspec={lr}}
    Categoria & Total \\
    P         & \num{300} \\
    Q         & \num{400} \\
  \end{tabelaabnt}
\end{tabela}
\endgroup

\end{document}
